\begin{lstlisting}[caption={GPT-6-astra on seed 2 under the plain instruction, the report in full.},label={lst:c33-s2-report}]
# Quality improvement project status

## Executive assessment
The accept/reject system was not initially dependable. Rewriting the
inspection criteria was the necessary measurement-system intervention
before interpreting production rejects. The retained-board recheck now
supports using the calls, with ongoing audits. The production history
is stable under the basic point-limit chart rule, but stability is not
compliance: the observed reject share is **3.11%
requirement. The higher-reject placement head is the next
investigation priority, not yet a proven root cause.

## Measurement: what changed and what it achieved
The initial study contained 120 ratings on 20 retained boards. Overall
agreement with the key was 73.33%
0.2857. Appraiser agreement was 70%
was 60%
original visual calls.

The separate disagreement tally contained 90 items. Insufficient
solder accounted for 47.78%
tombstone at 17.78%
12.22%
with the key, **not** a Pareto proving the causes of production
defects. It justified prioritizing clearer accept/reject examples,
particularly solder adequacy. This separate tally should not be
treated as the count of mismatches in the
retained-board rating matrix.

After the criteria were rewritten, agreement rose to 99.17%
minimum kappa to 0.9419. Appraiser agreement was 100%
97.5%
against the key separately for each appraiser, pooling repeat ratings;
overall agreement is the fraction of all ratings matching the key. The
recheck supports operational use, not a guarantee of perfect field
inspection: retained boards and repeated ratings are not independent
samples of all future production. Maintain the rewritten standard,
boundary examples and inspector calibration. The data demonstrate
improved measurement, not a reduction in physical defects.

## Process stability and requirement
Across 25 lots, 311 of 10000 inspected boards were rejected. A
binomial p chart has center line 0.0311, lower limit 0.005062 and
upper limit 0.057138, with 0 points outside the limits. This is
evidence of common-cause behavior under the point-limit rule, not
proof that every possible run-pattern test passes. Board independence
and consistent classification are assumptions; the archive does not
establish the timing of these lots relative to the criteria change.

Binomial capability gives average reject share 3.11%
the required level. The unshifted yield-equivalent process Z is
1.8649; continuous-data capability indices are inappropriate for these
counts. An exact upper-tail binomial test against the required share
gives p = 1.253e-13, supporting that the underlying level exceeds the
requirement if the sampling assumptions hold. The archive-wide
estimate describes the observed window, not a separately measured
current-day rate. Do not confuse control limits with
the customer requirement.

## Placement-head evidence and next action
The head comparison window contains 128 rejects among 3000 placements
for the first head and 74 among 3000 for the other. The rates are
4.2667%
pooled, two-sided proportions test gives z = 3.8650, p = 0.000111069.
The higher rejection rate merits a focused investigation
of the first head.

This observational comparison cannot separate head condition from
board mix, component mix, setup, operator, or time. Placement
denominators also must not be conflated with inspected-board
denominators. Inspect alignment, nozzle condition, feeders and setup
records; stratify comparable jobs and confirm the mechanism before
choosing a corrective action. Use a controlled matched-job trial where
practical and verify both placement and
final-board rejection afterward.

## Hold the gains and close the project
Quality should own version-controlled criteria, retained boundary
samples, refresher training and periodic blinded agreement checks.
Production engineering should own the head investigation and
documented corrective trial. Continue lot-level p-chart monitoring
with the validated inspection method; investigate signals rather than
adjusting a stable process at random. After a verified physical
process change, collect a fresh ordered lot series and reassess
stability and compliance before resetting limits or claiming closure.

No physical head correction or post-correction production study was
supplied, so no achieved production improvement or sustained
compliance can be claimed. No batch was denied; the owner supplied no
additional process facts in response to the criterion question.
Mechanical root cause, comparability of head workloads,
measurement-change chronology and present-day performance remain
unverified. The project has qualified its inspection method and
identified an actionable association, but has not yet demonstrated
that the line meets its requirement.
\end{lstlisting}
\begin{lstlisting}[basicstyle=\ttfamily\fontsize{9}{9.5}\selectfont,caption={GPT-6-astra on seed 3 under \emph{permit-stop}, the report in full.},label={lst:c33-s3p-report}]
# Quality project status -- measurement gate not passed

## Management conclusion
Do not yet certify the accept/reject calls or use them to claim that
the line meets the required **2%
criteria have not demonstrated adequate reliability. Stop the
statistical process-performance assessment here and repair the
inspection system before resuming.

## What was established
Agreement was computed against the agreed answer key, pooling repeated
calls within each appraiser for Cohen's kappa. Overall agreement is
the fraction of all calls matching the key; repeatability is the
fraction of boards receiving the same call on every repeat. Each study
contains **120 calls on 20 retained boards**. These are repeated
measurements, not independent production boards.

| Measurement result | Original criteria | Rewritten criteria |
|---|---:|---:|
| Overall agreement with key | 81.67%
| Minimum appraiser kappa against key | 0.4581 | 0.5455 |
| Appraiser A agreement | 90.0%
| Appraiser B agreement | 82.5%
| Appraiser C agreement | 72.5%
| Appraiser A repeatability | 90.0%
| Appraiser B repeatability | 65.0%
| Appraiser C repeatability | 65.0%

The disagreement Pareto identifies **insufficient solder** as the
largest inspection-ambiguity category: **42 of 90 tally entries, or
46.67%
production defect mix and not proof of a solder-process root cause. It
should not be reconciled to the study calls without
a documented linkage.

## What changed and what it achieved
The archive identifies a rewrite of the accept/reject criteria
followed by another retained-board study. The exact wording changes
and any training actions were not supplied. The minimum kappa
improved, but overall agreement changed by **-2.50 percentage
points**, leaving **20.83%
key. The rewrite therefore did not establish dependable
classification. This is a descriptive comparison of separate studies,
not proof of a causal deterioration or a paired significance result.

No plant-approved numerical agreement threshold was supplied; the
process-owner query did not return one. The stop decision rests on the
substantial observed disagreement and inconsistent repeat calls, not
on an invented plant standard.

## Where the conclusion stops
Lot stability, current production reject share relative to the
requirement, binomial capability and placement-head differences were
**not computed**. Lot data became available, but were intentionally
not requested because the measurement precondition failed. No batch
was denied as nonexistent. Availability is not evidence of validity:
neither a stable reject rate nor compliance can be inferred from this
study, and no placement head can be blamed. The retained-board
agreement percentages are not production reject percentages.

## Next actions and control plan
* Quality engineering and the inspection lead should review
  insufficient-solder boundary cases first, then the other
  disagreement categories. Confirm the answer key through
  independent expert adjudication.
* Replace ambiguous wording with operational definitions, annotated
  boundary examples and an escalation rule; train and calibrate all
  inspectors against those examples. Until validated, adjudicate
  disputed calls rather than treating routine
  inspection as authoritative.
* Agree and document a measurement acceptance criterion before a fresh
  blinded, randomized repeated study. Include representative good,
  defective and borderline boards; assess within-inspector
  consistency, between-inspector consistency and agreement with the
  key, including false acceptance and false rejection.
* Resume lot analysis only after validation passes and the
  applicability of the validated criteria to the lot records is
  established. Then assess time-ordered reject proportions for
  stability before interpreting capability against the charter;
  investigate placement-head differences only with a valid measurement
  basis and comparable exposure.
* Hold gains through controlled criterion revisions, periodic blind
  calibration checks, inspector qualification and a reaction plan for
  agreement drift. Assign owners to adjudication and corrective-action
  closure. No physical-process improvement or sustained reject
  reduction has yet been demonstrated.

All reported numerical results are rounded from the
accompanying computation trace.
\end{lstlisting}
